% !TeX root = ../papers/nc-p3.tex
\section{滑らかな平面三次点グラフと正規一次列の同値}
\label{sec:plane-cubic-point-graphs}
\begin{theorem}\label{thm:plane-cubic-point-graphs}
四辺の左右射影ファイバーが空か一点となる有限Frobenius窓で、点グラフが滑らかな平面三次曲線
$E$ とそれに交わらない有限スキームの和になるものは、次の全ての窓で尽きる。
次数3の偏極 $L$ を持つ $E$ と $e\in E\setminus E[3]$ に対し、平行移動 $T_e$ の
三生成Sklyanin商を取り、系\ref{cor:sklyanin-central-classification}の全ての中心一次PBW持上げを使う。
各窓は全次数で標準Koszul左右AS正則、裸周期1である。有限部は長さ8の三二次式の完全交叉で、
非被約でもよい。固定商ごとの添字保持同型類は同系の分裂点と
$\mathbb P(1,1,1,2,2,2)$ の有限群商であり、商を動かす同型は偏極曲線と平行移動の同型で判定する。
\end{theorem}
\begin{proof}[混合関係の両射影が全階数であること]
滑らかな種数1の三次埋込みは各位置で平面に入り、各平面を切る文字を $t_i$ とする。
その上の文字は次数3の偏極 $L_i$ の完全系である。隣接平面の積に制限した六関係の像は三次元。
実際、次元が3未満なら余分な曲面があり、次元は3を超えない。
後者は次数3の二完全系の積が全射であることによる。隣接偏極が等しければ平面対角を含むので等しくない。
異なる場合、前の偏極の基点自由鉛筆と後の偏極の積は鉛筆列と
$H^1(L_{i+1}/L_i)=0$ より全射となる。平面対積での三式の零スキームはこの三次グラフに等しい。
実際、曲面は存在せず、三式の完全交叉の両側次数は3で、既にこの曲線が全次数を占めている。

制限で零となる三関係を、平面上の座標 $X,Y$ を持ち上げて列で書けば
\[
          t_i B_iY-t_{i+1} A_iX+c_i t_it_{i+1}=0
\]
となる。$A_i,B_i$ は三行三列である。$A_i$ の核に射影点 $P$ があると仮定する。
$P$ が平面三次上なら、その次点との対で三混合式の微分の階数が1以下となり、
残る三式の微分を加えても六次元の周囲空間で階数4以下である。曲線はそこで滑らかなため矛盾する。
従って $P$ は平面三次の外にある。そこで混合三式は後ろの座標に関して恒等零、
残る三式は後ろの平面だけで階数3となる。従ってただ一つの次点が存在し、その点は平面の外にある。
点グラフの両射影の整合性から $P$ には前点も存在する。その前点も平面外でなければならない。
前の辺の混合式に代入すると $B_{i-1}P=0$ を得る。

この二つの階数欠損は許されない。$A_i$ の左核により $t_i\otimes v\in R_i$ なる非零な
純関係が生じ、$B_{i-1}$ の左核により $u\otimes t_i\in R_{i-1}$ が生じる。
従って三次交叉は純テンソル $u\otimes t_i\otimes v$ を含む。
しかし有限Frobenius対合により、これもある非零な第四文字評価による四次テンソルの縮約である。
各点評価で端の支持は少なくとも3である。隣接切断の像が全六関係であり、
その評価の階数が2以下ならば対応する射影ファイバーが直線を含むからである。
この矛盾で $A_i$ は可逆、逆向きで $B_i$ も可逆となる。

\smallskip
\noindent\textit{三次の単射性から一周期と全次数へ。}
混合式より $t_i V_{i+1}=V_i t_{i+1}$ は二次商の中で四次元となる。
実際、混合式は純な片側の関係を含まず、制限が零の関係は混合三式で尽きる。
よって $x t_{i+1}=t_i g_i(x)$ を与える可逆な $g_i:V_i\to V_{i+1}$ があり、
$g_i(t_i)=t_{i+1}$ である。両側から生成するイデアルは片側からも生成できる。
$B=A/(t_i)$ の三文字商は従って次元が少なくとも $20-10=10$ である。

$M_i=\ker(H^0(L_i)\otimes\mathcal O_E\to L_i)$ とおく。これは階数2、次数 $-3$ の安定束である。
次数 $-1$ 以上の直線部分束があれば、その自明三次元束への写像の成分は高々一つの切断の倍となる。
次数 $-1$ ならその共通零点で部分束でなく、次数0なら切断の一次独立性に反する。
$B$ の三次交叉は
\[
 H^0(M_i\otimes M_{i+2}\otimes L_{i+1})
 =\operatorname{Hom}(M_i^*,M_{i+2}\otimes L_{i+1})
\]
である。実際、三テンソルの中央の文字を曲線点に沿って評価すれば、
両隣接積が零という条件はちょうど両評価核束への包含になる（中央が完全切断系）。
右辺の二束は同じ傾きの安定束で、非零射は同型、次元は高々1である。
従って $\dim B_{i,i+3}=27-18+\dim H^0(M_i\otimes M_{i+2}\otimes L_{i+1})\le10$。
等号となり、$t$ の左右乗法は二次から三次へ単射である。
非零な商の三次交叉テンソルは両端の支持が3である。各点で上の束写像は同型なので、
その値の端支持は $M_i$ と $M_{i+2}$ のその点の二平面全体となり、これらは点を動かせば三空間を張る。

二文字関係の後ろに $t$ を付けて二回移動すると、今得た単射性で
$(g_i\otimes g_{i+1})R_i=R_{i+1}$ が従う。この一移動で基底を運べば窓は同じ通常関係を繰り返し、
$t$ は中心になる。商 $B$ は三次交叉の全支持と平面三次グラフから
補題\ref{lem:ternary-finite-graph-test}により三次元Koszul左右AS正則である。
商の関係 $r$ の持上げを $r=t\alpha(r)+t^2\beta(r)$ と書く。三次syzygyの二表示を引くと
定理\ref{thm:regular-normal-linear}の記法で $t(h+t b)=0$。
二次から三次への単射性で $h+t b=0$、商に移して $h\in R_B$ となる。
$h$ の持上げを代入し、$tV$ の四次元性を使えば $\alpha(h)=-b,\ \beta(h)=0$ となる。
従って同定理のPBW三条件が成立し、正則性は全次数に及ぶ。

\smallskip
\noindent\textit{商の自己同型と逆構成。}
中心表示の平面三次グラフの移動 $\rho$ は固定点を持たない。
固定点では中心三関係の二つの横方向に関する列が比例し、先と同じ接空間の矛盾が生じる。
楕円曲線で固定点を持たない自己同型は平行移動 $T_e$ であり、隣接偏極が異なるので $3e\ne0$ である。
逆にこの条件の $T_e$ に対し、完全切断系の三つの積関係はスキームとして滑らかな三次グラフを切る。
台については、積の汎関数の階数が1なら、命題\ref{prop:quintic-mixed-kernel}の拡大で
次数零の二階数束に二つ以上の独立切断がある。その行列式は $T_e^*L/L\ne1$ なので
二切断は同じ次数2以上の直線部分束に入り、汎関数は長さ1以下の因子上に支持される。
接空間については、$p$ で消える二平面の積が前後の偏極の積から $2p$ を引いた全切断を張る。
二つの次数2の完全鉛筆は相異なるため、鉛筆列で四次元から四次元への積が同型となるからである。
従ってこのグラフは被約である。

この三関係は通常の三生成Sklyanin形となる。実際、$E[3]$ は両偏極とグラフを保ち、
$H^0(L)$ 上で三指標の対角作用とその三つの巡回置換を与える。
これは位数3部分群での降下で各固有切断が一次元となり、他の三捩れがそれらを推移的に置換することによる。
全群が固有直線を保てば線型化により次数9の商に次数3の束が降下してしまう。
従って関係空間でも一つの対角指標の関係とその二つの巡回置換が基底である。
実際、文字の指標を一つ送れば二文字の指標は二つ送られ、三つが推移的になる。
三式は $A yz+B zy+C x^2$ とその巡回置換であり、通常の循環三次テンソルが全支持のsyzygyを与える。
よってこの商に系\ref{cor:sklyanin-central-classification}を適用できる。

全てのPBW持上げで、中心関係のため $t\ne0$ なら両点は等しく、
不等な平面三次の各点対の近くでは逆に両側で $t=0$ になる。
一方だけで $t=0$ の点対は存在しない。従って対角上の残余は射影三空間中の三つの二次式で、
無限遠では固定点がないため空、全体が有限である（完備な正次元部分は超平面に交わる）。
完全交叉で長さ8となる。これで全持上げの点グラフの仮定も確認した。
\end{proof}

\begin{corollary}
系\ref{cor:sklyanin-central-classification}の係数 $p,q,s$ と対称な一次補正行列 $M$、
定数列 $b$ を使うと、長さ8の孤立部はアフィン三空間における
\[
 \Phi(X)=\tfrac{s}{3}(X_1^3+X_2^3+X_3^3)+(p+q)X_1X_2X_3
           -\tfrac12 X^tMX-b^t X
\]
の臨界スキームである。一般の $b$ で八点は全て単純で、各特殊値の重複はこの勾配イデアルで決まる。
\end{corollary}
\begin{proof}
$t=1$ の対角へ三関係を制限すれば勾配に一致する。定数 $b$ を動かす勾配写像は
固有である。実際、そのグラフを射影三空間とパラメータ空間の積で閉じても
固定点のない最高次三式のため無限遠点がない。全ファイバーは長さ8で、有限かつ支配的、
標数零なので一般点でエタールである。
\end{proof}
